\documentclass[11pt]{article}
\usepackage[a4paper,margin=25mm]{geometry}
\usepackage{amsmath,amssymb,amsthm}
\usepackage[hidelinks]{hyperref}
\newtheorem{theorem}{Theorem}
\newtheorem{lemma}{Lemma}
\title{Markov Numbers Do Not Have Square-Root Counting Growth:\\A Disproof of Conjecture 00000008333}
\author{GitHub: gaochengzhecpu}
\date{}
\begin{document}
\maketitle
\begin{abstract}
Let $M(X)$ count the distinct positive integer Markov numbers at most
$X$. We prove the explicit estimate
\[
 M(64^n)\le 2+3(6n+1)^2\,2^{2n+1}\qquad(n\ge0).
\]
Consequently $M(64^n)/8^n\longrightarrow0$, contradicting
$M(X)\sim c\sqrt X$ for every $c>0$. The argument uses elementary
descent, a complete binary tree of actual Markov triples, and a bound
on its words after their second right branch. Lean verifies the full
bridge from positive integer solutions to the actual counting function
and the quantified negation of the asserted asymptotic.
\end{abstract}

\section*{Reading of the source}
The English source asserts square-root growth for the distribution
of Markov numbers; the Chinese version explicitly writes
$\mathrm{mark\_dist}\sim c x^{1/2}$. We use the usual cumulative
counting interpretation
\[
\begin{split}
\mathcal M&=\{m\in\mathbb Z_{>0}:\ \exists a,b\in\mathbb Z_{>0},\
 a^2+b^2+m^2=3abm\},\\
M(X)&=\#\{m\in\mathcal M:m\le X\}.
\end{split}
\]
Numbers are counted once, even if they occur in several triples. The
constant in an asymptotic growth law is positive. It suffices to
consider integer cutoffs: an asserted asymptotic for all real cutoffs
would also hold along the integers.

The source also discusses prime Markov numbers, Frobenius uniqueness,
and random walks. We refute its square-root clause, which is necessary
for the conjunction asserted there. We make no claim to settle its
other clauses. In particular, the proof neither assumes nor establishes
uniqueness of the largest coordinate among tree nodes.

For historical context, Zagier proved a much sharper asymptotic for
the number of sorted Markov triples below a bound~\cite{zagier}.
That count allows repeated largest-coordinate labels. We do not
identify it with the count of distinct numbers or use that asymptotic
in this proof. The elementary upper bound below is verified directly.

\section*{Descent and complete tree coverage}
Call $(a,b,c)$ a sorted Markov triple when $0<a\le b\le c$ and
\begin{equation}\label{eq:markov}
 a^2+b^2+c^2=3abc.
\end{equation}
Starting with $T_\varnothing=(1,2,5)$, define children of a strict
sorted triple by
\[
 L(a,b,c)=(a,c,3ac-b),\qquad
 R(a,b,c)=(b,c,3bc-a).
\]
For a finite word $w$ in $L,R$, let $T_w$ be its value, applying its
letters from left to right, and let $C(w)$ be its largest coordinate.

\begin{lemma}\label{lem:coverage}
Every sorted Markov triple is $(1,1,1)$, $(1,1,2)$, or some $T_w$.
Every $T_w$ is strict: $0<a<b<c$.
\end{lemma}
\begin{proof}
The substitution $b\mapsto3ac-b$ preserves \eqref{eq:markov}; this
follows by expansion. For a strict triple, $a\ge1$ and $b<c$, so
$3ac-b>c$. The same argument gives $3bc-a>c$. Thus both children
are strict Markov triples, and the root is one as well.

For completeness, take a sorted solution other than one with
$a=b=1$, and set $d=3ab-c$. First $c<3ab$: the opposite inequality
would give $3abc\le c^2$, contradicting \eqref{eq:markov} and
$a,b>0$. Hence $d>0$, and
\[
 d+c=3ab,\qquad dc=a^2+b^2.
\]
We claim $d<b$. Since $(a,b)\ne(1,1)$,
\begin{equation}\label{eq:negative}
 a^2+2b^2<3ab^2.
\end{equation}
Indeed, for $a=1$ we have $b\ge2$; for $a\ge2$ we use
$a^2\le b^2$ and $3ab^2\ge6b^2>3b^2$. If $d\ge b$, then
$(d-b)(c-b)\ge0$, whereas the displayed identities give
\[
 (d-b)(c-b)=a^2+2b^2-3ab^2<0.
\]
This is a contradiction. Also $b<c$, because $b=c$ would imply
$db=a^2+b^2>b^2$ and hence $d>b$.

Mutation shows that $(a,b,d)$ again satisfies the equation. If
$d\le a$, its sorted version is $(d,a,b)$, and the original triple
is its $R$ child. If $d>a$, its sorted version is $(a,d,b)$, and the
original is its $L$ child. Either parent has smaller coordinate sum.
Induction therefore reduces to the case $a=b=1$, where the equation
gives $c=1$ or $2$. The parent $(1,1,1)$ can lead only to $(1,1,2)$;
the first strict triple after $(1,1,2)$ is $(1,2,5)$.
These are exactly the exceptional triples and root stated above.
\end{proof}

\begin{lemma}\label{lem:birth}
For every $m\in\mathcal M$, either $m=1$, $m=2$, or $C(w)=m$ for
some finite word $w$.
\end{lemma}
\begin{proof}
Sort any positive solution containing $m$ and use
Lemma~\ref{lem:coverage}. At the root, its coordinates are $1,2,5$,
and $5=C(\varnothing)$. Each child reuses two coordinates from its
parent and creates just its own new largest coordinate. Induction
along its path therefore shows that every coordinate is either
$1$, $2$, or the largest coordinate of an ancestor. This argument
does not bound the largest coordinate of an arbitrary triple
containing a small $m$: it chooses the ancestor at which $m$ is a
largest coordinate.
\end{proof}

\section*{A bound on words below a cutoff}
For a strict triple $(a,b,c)$, either child's largest coordinate is at
least $(3a-1)c$. For $L$ this follows from $b\le c$; for $R$ it
follows from $b\ge a$ and $a\le c$. A child never decreases the
least coordinate. In addition, every $R$ increases that coordinate
by at least one, since it replaces $a$ by $b>a$.
It follows that
\begin{enumerate}
\item $C(w)\ge 2^{|w|}\cdot5$, so $C(w)\le64^n$ implies
$|w|\le6n$;
\item after two $R$ letters the least coordinate is at least three,
so each remaining step increases the largest coordinate by a factor
of at least eight. If $w=uv$, $u$ contains at least two $R$ letters,
and $C(w)\le64^n$, then $8^{|v|}\le64^n$ and $|v|\le2n$.
\end{enumerate}
These statements follow equally for any initial strict triple by
induction, multiplying its largest coordinate by the stated factor
at each step.

Write $A=6n+1$. Every word with $C(w)\le64^n$ belongs to one of
three classes, according to its first two $R$ letters:
\[
 L^i,\qquad L^i R L^j,\qquad L^i R L^j Rv.
\]
Here $0\le i,j\le6n$, and in the third class $|v|\le2n$.
The first two classes explicitly include paths stopping before the
second $R$. There are at most $A$ choices of the first type,
$A^2$ of the second, and
\[
 A^2\sum_{k=0}^{2n}2^k\le A^2\,2^{2n+1}
\]
of the third. Thus the total is at most
$3A^2\,2^{2n+1}$. We need only an upper bound: different words may
have the same largest coordinate.

By Lemma~\ref{lem:birth}, every distinct $m\in\mathcal M$ at most
$64^n$, except possibly $1,2$, is the largest coordinate of a word
in this finite collection. Taking its image under $C$ gives
\begin{equation}\label{eq:bound}
 M(64^n)\le2+3(6n+1)^2\,2^{2n+1}.
\end{equation}
This is a bound on the actual set of positive integer Markov numbers,
not just a count of an independently specified tree.

\section*{Contradiction to square-root growth}
Dividing \eqref{eq:bound} by $8^n=\sqrt{64^n}$ gives
\[
0\le\frac{M(64^n)}{8^n}\le
\frac{2}{8^n}+216\frac{n^2}{2^n}
+72\frac n{2^n}+6\frac1{2^n}\longrightarrow0.
\]
Each term tends to zero because a fixed polynomial grows more slowly
than a geometric sequence of ratio greater than one. In particular,
for every $c>0$ and every $N$ there is $n\ge N$ with
$M(64^n)<c8^n$.

\begin{theorem}
There is no real $c>0$ for which $M(X)/\sqrt X\longrightarrow c$
as the integer cutoff $X\longrightarrow\infty$.
\end{theorem}
\begin{proof}
Such a limit, restricted to $X=64^n\longrightarrow\infty$, would
give $M(64^n)/8^n\longrightarrow c$. The same sequence tends to
zero, so uniqueness of limits would give $c=0$, a contradiction.
Thus the square-root clause, and hence the source's conjunction,
is false under the stated counting interpretation.
\end{proof}

\section*{Formal verification and scope}
The Lean development defines the actual equation over natural
numbers with all three coordinates positive. Its finite set at
cutoff $X$ filters $\{0,\ldots,X\}$ by existence of positive
$a,b$ satisfying the equation; its cardinal is exactly $M(X)$.
No table of Markov numbers is used. The descent, sorted tree
coverage, six sorting cases, coordinate birth lemma, path growth,
three word classes, finite-set image bound, real limit, and final
negation are all formally proved.

The principal statements, in namespace \texttt{Markov8333}, are
\begin{center}
\texttt{number\_birth}, \texttt{count\_bound},\\
\texttt{count\_subsequence\_zero}, \texttt{arbitrarily\_small\_ratio},\\
\texttt{conjecture\_8333\_false :}\quad
$\neg\,\texttt{SquareRootGrowth}$.
\end{center}
Here \texttt{SquareRootGrowth} is the existential positive-constant
limit in the theorem above, with the actual finite-set cardinal.
Lean also proves $\sqrt{64^n}=8^n$. The final theorem is not a
conditional statement assuming completeness of the tree or
uniqueness of Markov labels. The interpretation of the source's
word ``distribution'' as a cumulative count, and the observation
that a real-cutoff asymptotic implies its integer restriction,
are stated explicitly here. No assertion about the other three
source clauses is formalized or needed.

The Lean source contains no unproved placeholders or additional axioms.
The reported axioms of the five audited theorems are only
\texttt{propext}, \texttt{Classical.choice}, and \texttt{Quot.sound}.
Development and self-review were performed by a delegated AI agent;
the coordinating agent reviews separately. No external independent
review is claimed.

\section*{Source and reproduction}
The exact bilingual statement is archived as \texttt{SOURCE.md},
from the competition repository at commit
\begin{center}\texttt{b355e42906e6eaf52bc7e9d005e0ef02d8aa7400}.\end{center}
The source path is
\begin{center}\texttt{conjectures/00000008333.md}.\end{center}
Source hashes and snapshot checks are recorded in the accompanying
provenance file. No earlier submission was found in that snapshot.

The project pins Lean 4.19.0 and Mathlib commit
\begin{center}\texttt{c44e0c8ee63ca166450922a373c7409c5d26b00b}.\end{center}
From the supplied \texttt{lean} directory, run
\begin{center}\texttt{lake env lean -DwarningAsError=true Main.lean}.\end{center}
The proof is self-contained apart from standard finite-set,
arithmetic, and real-analysis facts supplied by Mathlib.
\begin{thebibliography}{1}
\bibitem{zagier} D. Zagier, \emph{On the Number of Markoff Numbers
Below a Given Bound}, Mathematics of Computation 39 (160), 1982,
709--723.
\href{https://people.mpim-bonn.mpg.de/zagier/files/doi/10.2307/2007348/fulltext.pdf}{Author's public copy}.
\end{thebibliography}
\end{document}
